\section*{Chapter \ref{ch:g6:fractions} --- Fractions: First Steps}

\begin{solution}{exo:g6:fractions:1}
Cake: $\frac38$ eaten. \quad Chocolate: $\frac{7}{12}$. \quad
Pie: $\frac66 = 1$ (the whole pie).
\end{solution}

\begin{solution}{exo:g6:fractions:2}
Four parts out of five are colored. $\frac45 < 1$ (four parts out of
five is less than the whole bar); $\frac65 > 1$ (six fifths is one
whole bar and one fifth more).
\end{solution}

\begin{solution}{exo:g6:fractions:3}
On the line graduated in thirds: $\frac13$ is one tick after $0$;
$\frac53$ is one tick before $2$ (since $2 = \frac63$); $\frac73$ is
one tick after $2$. So $\frac53$ and $\frac73$ are both at distance
$\frac13$ from $2$.
\end{solution}

\begin{solution}{exo:g6:fractions:4}
$\frac94 = 2 + \frac14$; \quad
$\frac{13}{5} = 2 + \frac35$; \quad
$\frac{12}{6} = 2$ (exactly); \quad
$\frac{25}{8} = 3 + \frac18$.
\end{solution}

\begin{solution}{exo:g6:fractions:5}
$\frac23 = \frac{8}{12}$ (multiply by $4$); \quad
$\frac{15}{20} = \frac34$ (divide by $5$); \quad
$\frac{7}{10} = \frac{21}{30}$ (multiply by $3$); \quad
$\frac56 = \frac{10}{12}$ (divide by $2$).
\end{solution}

\begin{solution}{exo:g6:fractions:6}
$\frac68 = \frac34$; \quad
$\frac{15}{25} = \frac35$; \quad
$\frac{18}{24} = \frac{9}{12} = \frac34$; \quad
$\frac{35}{7} = 5$ (a whole number).
\end{solution}

\begin{solution}{exo:g6:fractions:7}
$\frac12$ of $86$: $86 \div 2 = 43$.

$\frac34$ of $60$: $60 \div 4 = 15$, then $15 \times 3 = 45$.

$\frac25$ of $35$: $35 \div 5 = 7$, then $7 \times 2 = 14$.

$\frac{7}{10}$ of $250$: $250 \div 10 = 25$, then $25 \times 7 = 175$.
\end{solution}

\begin{solution}{exo:g6:fractions:8}
$\frac12 > \frac13$: halves are bigger parts than thirds.
$\frac35 > \frac25$: same parts, more of them.
$\frac34 = \frac68 < \frac78$: seven eighths beat six eighths.
\end{solution}

\begin{solution}{exo:g6:fractions:9}
$\frac12 = 0.5$; \quad $\frac34 = 0.75$; \quad
$\frac{7}{10} = 0.7$; \quad $\frac94 = 2.25$. The \omterm{def:g3:numbers:evenodd}{odd} one out is
$\frac13 = 0.333\dots$, whose decimal writing never ends: it is not a
\omterm{def:g6:decimals:places}{decimal number}.
\end{solution}

\begin{solution}{exo:g6:fractions:10}
Bus users: $\frac34$ of $28$: $28 \div 4 = 7$, then
$7 \times 3 = 21$ students. Walkers: $28 - 21 = 7$ students (that is
the remaining \omterm{def:g3:division:half}{quarter}: $28 \div 4 = 7$). Check: $21 + 7 = 28$.
\end{solution}

\begin{solution}{exo:g6:fractions:11}
Two thirds of the savings are $18$ euros, so \emph{one} third is
$18 \div 2 = 9$ euros, and the whole savings (three thirds) were
$9 \times 3 = 27$ euros. Check: $\frac23$ of $27$ is $18$.
\end{solution}

\begin{solution}{exo:g6:fractions:12}
The used water went from $\frac58$ to $\frac48$ (a \omterm{def:g3:division:half}{half} is four
eighths) of the tank: the $10$ liters represent
$\frac58 - \frac48 = \frac18$ of the \omterm{def:g2:measure:units}{capacity}. So the full tank holds
$10 \times 8 = 80$ liters. Check: $\frac58$ of $80$ is $50$, minus
$10$ leaves $40$, which is indeed \omterm{def:g3:division:half}{half} of $80$.
\end{solution}

\begin{solution}{pb:g6:fractions:1}
\textbf{1.} \omterm{def:g3:division:half}{Half} of $16$ squares is $8$ squares. As a \omterm{def:g6:fractions:def}{fraction}:
$\frac12$ of the bar, and since each \omterm{def:g3:division:half}{half} is $8$ squares out of
$16$, also $\frac{8}{16}$ --- the same \omterm{def:g6:fractions:def}{fraction} written two ways
(\cref{prop:g6:fractions:equal}).

\textbf{2.} \omterm{def:g3:division:half}{Half} of the $8$ remaining squares is $4$ squares,
which is $\frac{4}{16}$ of the bar. Simplifying by $4$:
$\frac{4}{16} = \frac14$. On the picture: cutting the remaining
\omterm{def:g3:division:half}{half} in two produces two \omterm{def:g3:division:half}{quarters} of the original bar --- \omterm{def:g3:division:half}{half} of
a \omterm{def:g3:division:half}{half} is a \omterm{def:g3:division:half}{quarter}.

\textbf{3.} Third bite: \omterm{def:g3:division:half}{half} of $4$ squares is $2$ squares, i.e.\
$\frac{2}{16} = \frac18$ of the bar. Fourth bite: $1$ square,
i.e.\ $\frac{1}{16}$. From one bite to the next, the \omterm{def:g4:fractions:def}{denominator}
\emph{doubles}: $2$, $4$, $8$, $16$.

\textbf{4.} Eaten: $8 + 4 + 2 + 1 = 15$ squares out of $16$, that
is $\frac{15}{16}$ of the bar. Left: $1$ square, $\frac{1}{16}$.

\textbf{5.} $\frac12 = \frac{8}{16}$ (multiply top and bottom by
$8$), $\frac14 = \frac{4}{16}$ (by $4$), $\frac18 = \frac{2}{16}$
(by $2$), and $\frac{1}{16}$ stays. Counting sixteenths:
$8 + 4 + 2 + 1 = 15$ of them, so the total is $\frac{15}{16}$.

\textbf{6.} Bites: $32$, $16$, $8$, $4$, $2$, $1$ squares. Total
eaten: $32 + 16 + 8 + 4 + 2 + 1 = 63$ squares, so $1$ square
remains: $\frac{1}{64}$ of the bar.

\textbf{7.} After each bite there remains
\[
\frac12,\quad \frac14,\quad \frac18,\quad \frac{1}{16},\quad
\frac{1}{32},\quad \frac{1}{64} :
\]
the \omterm{def:g3:division:remainder}{remainder} halves each time --- its \omterm{def:g4:fractions:def}{denominator} doubles. A
seventh, imaginary bite would leave $\frac{1}{128}$.

\textbf{8.} Every bite eats only \emph{\omterm{def:g3:division:half}{half}} of what remains, so
the other \omterm{def:g3:division:half}{half} of the \omterm{def:g3:division:remainder}{remainder} is still there: after any number
of bites, something is always left. The eaten total is therefore
always below $1$ --- exactly like the twenty nines of
\cref{pb:g6:decimals:1}, which come ever closer to $3$ without
reaching it.

\textbf{9.} Both \omterm{def:g6:fractions:def}{fractions} have \omterm{def:g4:fractions:def}{numerator} $1$, and cutting into
$128$ parts makes smaller parts than cutting into $100$: so
$\frac{1}{128} < \frac{1}{100}$. After the seventh bite the
\omterm{def:g3:division:remainder}{remainder} is $\frac{1}{128}$, so the eaten part exceeds
$1 - \frac{1}{100} = \frac{99}{100}$: the challenge is won at
bite seven.

\textbf{10.} In sixteenths, the totals are $\frac{8}{16}$,
$\frac{12}{16}$, $\frac{14}{16}$, $\frac{15}{16}$. Each new total
lands exactly \emph{halfway} between the previous total and $1$:
the staircase keeps halving its remaining distance to $1$,
without ever stepping on it.

\textbf{11.} What is right: at every stage of Zeno's description,
some distance indeed remains (question 8) --- the list of stages
never ends. The trap: the stages become extremely short, in
distance \emph{and} in running time; the runner does not spend
equal time on each stage. Describing the run in infinitely many
shrinking pieces does not make the run itself endless: the
runner reaches the wall, and question 9 shows the description
passing any point short of it.

\textbf{12.} Cut $8$ bars into halves: $16$ half-pieces. Cut $4$
bars into \omterm{def:g3:division:half}{quarters}: $16$ quarter-pieces. Cut $2$ bars into
eighths: $16$ eighth-pieces. Cut the last bar into sixteenths:
$16$ small pieces. That uses $8 + 4 + 2 + 1 = 15$ bars, and each
child receives one piece of each size.

\textbf{13.} Each child's share is
\[
\frac12 + \frac14 + \frac18 + \frac{1}{16} = \frac{15}{16}
\]
of a bar (question 5). Sixteen such shares make
$16 \times \frac{15}{16} = 15$ bars
(\cref{thm:g6:fractions:quotient}): exactly what was cut, with
nothing left over --- a fair sharing of $15$ bars among $16$
children, by halving alone.

\textbf{14.} The neighbours receive $\frac78$ of a bar each. With
\omterm{def:g4:fractions:def}{denominator} $16$: $\frac78 = \frac{14}{16}$, while our children
receive $\frac{15}{16}$. So a child of the $16$ receives more:
$\frac{15}{16} > \frac{14}{16}$, by one sixteenth of a bar.

\textbf{15.} In the square picture, each shaded bite fills \omterm{def:g3:division:half}{half}
of the still-unshaded corner, and the unshaded corner keeps
shrinking: after each bite it is \omterm{def:g3:division:half}{half} as large as before, and it
never disappears. The two true facts: (i) after every finite
number of bites, a \omterm{def:g3:division:remainder}{remainder} is left --- the total eaten is
always strictly below $1$; (ii) that \omterm{def:g3:division:remainder}{remainder} becomes smaller
than any \omterm{def:g6:fractions:def}{fraction} one cares to name ($\frac{1}{100}$,
$\frac{1}{1000}$, \dots), if one bites long enough. Giving the
words ``the infinite \omterm{def:g1:addition:def}{sum} equals $1$'' their exact meaning is the
job of \emph{limits}, in the High School volume.
\end{solution}
